\chapter{What to Do: Antifragile Strategies Condizionate ai Quattro
Mondi}\label{cosa-fare-strategie-antifragili-condizionate-ai-quattro-mondi}

\begin{quote}
\textbf{Core argument:} la DEQ$^2$ non termina nella diagnosi. Lo
spazio degli scenari M1-M4 del Ch.~10.4 è anche uno spazio di
\emph{strategie ammissibili} --- diverse per ciascun livello
(individuale, organizzativo, statale) e per ciascun orizzonte (tattico,
operational, strategico). Questo chapter organizza ventisette protocolli
operativi (3 livelli $\times$ 3 orizzonti $\times$ 3 scenari aggregati significativi)
come \emph{traiettorie programmate} nel sottospazio
\(\mathbb{V}^{(4)}_{\text{ammissibile}}\), distinguendo strategie
\emph{passive} (massimizzano payoff atteso entro lo scenario realizzato)
da strategie \emph{antifragili} (massimizzano payoff atteso
\emph{through} lo spazio degli scenari, sfruttando l'asimmetria di
Taleb).
\end{quote}



\section{Strategie Passive vs Antifragili nel Linguaggio
DEQ$^2$}\label{strategie-passive-vs-antifragili-nel-linguaggio-dequxb2}

\Hypoth{} \textbf{Definizione operativa.} Una strategia
\(\boldsymbol{\sigma}\) applicata a un punto
\(\boldsymbol{p}(t) \in \mathbb{V}^{(4)}\) produce una traiettoria
\(\boldsymbol{p}_{\boldsymbol{\sigma}}(t+\Delta t)\) e un payoff atteso
\(\mathbb{E}[\Pi_{\boldsymbol{\sigma}}]\). Distinguiamo:

\begin{itemize}
\tightlist
\item
  \textbf{Strategia passiva}: ottimizza
  \(\mathbb{E}[\Pi_{\boldsymbol{\sigma}} \mid \mathcal{M}_k]\) assumendo
  dato lo scenario \(\mathcal{M}_k \in \{M_1, M_2, M_3, M_4\}\);
\item
  \textbf{Strategia antifragile}: ottimizza
  \(\mathbb{E}_{\mathcal{M}}[\Pi_{\boldsymbol{\sigma}}] = \sum_k \Pr(\mathcal{M}_k) \cdot \mathbb{E}[\Pi_{\boldsymbol{\sigma}} \mid \mathcal{M}_k]\)
  con vincolo
  \(\partial^2 \Pi_{\boldsymbol{\sigma}} / \partial \sigma^2 > 0\)
  (convessità with respect tolla volatilità sistemica).
\end{itemize}

\Proved{} \textbf{Risultato structural.} Una strategia antifragile è
\emph{quasi sempre} sub-ottimale in ciascun scenario specifico --- il
suo vantaggio emerges solo quando \emph{lo scenario realizzato è
incerto}. Per il sistema 2026-2031 con
\(\Pr(\text{decoerenza globale}) \approx 0{,}80\)-\(0{,}85\) ma
\(\Pr(\Phi\text{-engine}) \approx 0{,}55\)-\(0{,}70\), l'incertezza
sullo scenario realizzato è massima fra le quattro coppie di mondi,
therefore il valore relativo dell'antifragility è massimo \emph{adesso}.

\Hypoth{} \textbf{Confollowsnza.} I ventisette protocolli del Ch.~12 sono
presentati come \emph{strategie antifragili by design}: ciascuno è
progettato per produrre payoff non-negativo in \emph{almeno tre dei
quattro mondi} M1-M4. La verification della condizione è esplicitata in
table per ciascun protocollo.



\section{La Matrice Tre-Tre: Livelli per
Orizzonti}\label{la-matrice-tre-tre-livelli-per-orizzonti}

\Hypoth{} \textbf{Costruzione.} I tre livelli e i tre orizzonti
definiscono nove \emph{celle} di intervento. Ciascuna cella riceve un
protocollo generico (replicabile) più tre specializzazioni ai mondi (M1,
M2, M3 --- il caso M4 è omesso because lo scenario è il più stabile e
richiede strategie meno specifiche).

{\def\LTcaptype{none} % do not increment counter
{\small\begin{longtable}[]{@{}
  >{\raggedright\arraybackslash}p{(\linewidth - 6\tabcolsep) * \real{0.2500}}
  >{\raggedright\arraybackslash}p{(\linewidth - 6\tabcolsep) * \real{0.2500}}
  >{\raggedright\arraybackslash}p{(\linewidth - 6\tabcolsep) * \real{0.2500}}
  >{\raggedright\arraybackslash}p{(\linewidth - 6\tabcolsep) * \real{0.2500}}@{}}
\toprule\noalign{}
\begin{minipage}[b]{\linewidth}\raggedright
\end{minipage} & \begin{minipage}[b]{\linewidth}\raggedright
\textbf{Tattico} (mesi)
\end{minipage} & \begin{minipage}[b]{\linewidth}\raggedright
\textbf{Operativo} (1-3 anni)
\end{minipage} & \begin{minipage}[b]{\linewidth}\raggedright
\textbf{Strategico} (3-15 anni)
\end{minipage} \\
\midrule\noalign{}
\endhead
\bottomrule\noalign{}
\endlastfoot
\textbf{Individuale} & I-T: opzionalità minima & I-O: portafoglio
multi-vettoriale & I-S: traiettoria di vita anti-eteronomica \\
\textbf{Organizzativo} & O-T: redistribuzione decisionale & O-O:
ridondanza modulare & O-S: missione orientata ai mondi-target \\
\textbf{Statale} & S-T: protocollo di coerenza in crisi & S-O: politica
estera multi-vettoriale & S-S: investimento structural in
\(\Phi\)-engine \\
\end{longtable}}
}

\Proved{} \textbf{Lettura.} Ciascuna cella corrisponde a un
\emph{intervento mirato} su una o più variables di \(\mathbb{V}^{(4)}\).
Per example I-S agisce su \(A^{\text{individuale}}\) (resilience
personale) e su \(\varepsilon_{\text{dis}}^{\text{individuale}}\)
(resistenza al disimpegno). S-O agisce su \(\Phi^{\text{esportabile}}\)
(capacità di mediazione) e su \(\Omega^{\text{nazionale}}\) (riduzione
conflitti improduttivi). S-S agisce su \(\mathcal{S}\) stessa
(modificazione dei parameters esogeni \(s_{iii}, z, \delta\)).



\section{Strategie Individuali (livello
I)}\label{strategie-individuali-livello-i}

\subsection{I-T --- Opzionalità Minima (orizzonte
mesi)}\label{i-t-opzionalituxe0-minima-orizzonte-mesi}

\Hypoth{} \textbf{Protocollo.} Mantenere, in ogni dato istante, almeno
tre opzioni concrete e attivabili in \(< 90\) giorni --- relative a
\emph{una} delle followsnti categorie: residenza, professione, valuta di
risparmio, rete di affidamento personale.

{\def\LTcaptype{none} % do not increment counter
{\small\begin{longtable}[]{@{}
  >{\raggedright\arraybackslash}p{(\linewidth - 2\tabcolsep) * \real{0.5000}}
  >{\raggedright\arraybackslash}p{(\linewidth - 2\tabcolsep) * \real{0.5000}}@{}}
\toprule\noalign{}
\begin{minipage}[b]{\linewidth}\raggedright
Mondo
\end{minipage} & \begin{minipage}[b]{\linewidth}\raggedright
Specializzazione
\end{minipage} \\
\midrule\noalign{}
\endhead
\bottomrule\noalign{}
\endlastfoot
M1 (decoherence + $\Phi$-engine) & priorità a residenza/cittadinanza in attori
\(\mathcal{R}_+\) (Brazil, alcuni paesi UE non-frontali) \\
M2 (decoherence senza $\Phi$-engine) & priorità a valuta di risparmio
diversificata (CHF, oro, basket BRICS) \\
M3 (rottura Chinese precoce) & priorità a rete di affidamento personale
fuori da catene di fornitura cinesi \\
\end{longtable}}
}

\Proved{} \textbf{Condizione antifragile.} L'opzionalità minima è
positiva in M1, M2, M3, M4 --- \emph{sempre}. Costo: \(5\)-\(15\%\) del
payoff in ciascun scenario specifico (overhead di mantenimento delle
opzioni). Beneficio: \(> 50\%\) del payoff in caso di scenario
\emph{non} anticipato.

\subsection{I-O --- Portafoglio Multi-Vettoriale (orizzonte 1-3
anni)}\label{i-o-portafoglio-multi-vettoriale-orizzonte-1-3-anni}

\Hypoth{} \textbf{Protocollo.} Strutturare attività professionale,
finanziaria e relazionale come \emph{portafoglio} di vettori
parzialmente decorrelati: nessun vettore \(> 35\%\) del totale, almeno
tre vettori a correlazione \(< 0{,}3\) fra loro.

{\def\LTcaptype{none} % do not increment counter
{\small\begin{longtable}[]{@{}
  >{\raggedright\arraybackslash}p{(\linewidth - 2\tabcolsep) * \real{0.5000}}
  >{\raggedright\arraybackslash}p{(\linewidth - 2\tabcolsep) * \real{0.5000}}@{}}
\toprule\noalign{}
\begin{minipage}[b]{\linewidth}\raggedright
Mondo
\end{minipage} & \begin{minipage}[b]{\linewidth}\raggedright
Specializzazione
\end{minipage} \\
\midrule\noalign{}
\endhead
\bottomrule\noalign{}
\endlastfoot
M1 & un vettore \(\Phi\)-engine (es. economia Brazilian, BRICS-related,
transizione climatica), uno tech-sovrano, uno tradizionale \\
M2 & priorità a vettori non-correlati ai cicli di decoherence (beni
reali, comunità local, capacità manuale) \\
M3 & priorità a vettori indipendenti dalla supply chain Chinese
(manufacturing local, AI europea/USA/Brazilian) \\
\end{longtable}}
}

\Proved{} \textbf{Condizione antifragile.} Un portafoglio di tre vettori
a correlazione \(< 0{,}3\) ha varianza \(\sim\) varianza media \(/ 3\);
in regime di alta volatilità sistemica (\(\sigma \uparrow\)), la
riduzione di varianza individuale è strutturalmente positiva (Markowitz
applicato a vita personale).

\subsection{I-S --- Traiettoria di Vita Anti-Eteronomica (orizzonte 3-15
anni)}\label{i-s-traiettoria-di-vita-anti-eteronomica-orizzonte-3-15-anni}

\Hypoth{} \textbf{Protocollo.} Costruire competenze, relazioni e
riferimenti culturali tali da rendere \emph{non sostituibile} il proprio
progetto di vita with respect to una traiettoria di disimpegno morale ---
operativizzazione individuale della \emph{critical resistance} \(R_c\)
del Ch.~9.4.

{\def\LTcaptype{none} % do not increment counter
{\small\begin{longtable}[]{@{}
  >{\raggedright\arraybackslash}p{(\linewidth - 2\tabcolsep) * \real{0.5000}}
  >{\raggedright\arraybackslash}p{(\linewidth - 2\tabcolsep) * \real{0.5000}}@{}}
\toprule\noalign{}
\begin{minipage}[b]{\linewidth}\raggedright
Componente
\end{minipage} & \begin{minipage}[b]{\linewidth}\raggedright
Indicatore
\end{minipage} \\
\midrule\noalign{}
\endhead
\bottomrule\noalign{}
\endlastfoot
Competenza non sostituibile da AI generale entro 10 anni & almeno una
abilità con curva di apprendimento \(> 5\) anni \\
Rete di relazioni significative non transazionali & almeno cinque
relazioni decennali sostenute oltre l'utilità \\
Reference culturale non riducibile a tendenza & radicamento esplicito
in tradizione (filosofica, religiosa, artistica) di almeno cento anni \\
\end{longtable}}
}

\Proved{} \textbf{Condizione antifragile.} L'identità anti-eteronomica
produce \(\varepsilon_{\text{dis}}^{\text{individuale}} \to 0\) stabile
--- proprietà che ha valore \emph{positivo} in tutti i mondi M1-M4 e
\emph{cresce} in valore quando
\(\varepsilon_{\text{dis}}^{\text{sistemico}} \uparrow\).



\section{Strategie Organizzative (livello
O)}\label{strategie-organizzative-livello-o}

\subsection{O-T --- Redistribuzione Decisionale in Crisi (orizzonte
mesi)}\label{o-t-redistribuzione-decisionale-in-crisi-orizzonte-mesi}

\Hypoth{} \textbf{Protocollo.} In presenza di pattern WE-DEQ$^2$ (Cap.
10.2), redistribuire entro \(\leq 60\) giorni l'autorità decisionale
dall'apice verso almeno tre nodi periferici --- \emph{non} per delega ma
per \emph{autonomia conditional} (``se accade \(X\), decidi tu senza
reference centrale entro \(Y\) giorni'').

\Proved{} \textbf{Condizione antifragile.} Aumenta
\(\Phi^{\text{interno}}\) dell'organizzazione through riduzione di
\(\zeta(t)\) esogeno (ogni nodo periferico assorbe shock localmente) e
aumento di \(A^{\text{organizzativo}}\) (più punti di risposta
autonoma).

\subsection{O-O --- Ridondanza Modulare (orizzonte 1-3
anni)}\label{o-o-ridondanza-modulare-orizzonte-1-3-anni}

\Hypoth{} \textbf{Protocollo.} Per ciascuna funzione critica
dell'organizzazione, mantenere \(\geq 2\) implementazioni indipendenti,
una preferibilmente in tecnologia/giurisdizione/cultura
\emph{non-dominante} nello scenario realizzato.

{\def\LTcaptype{none} % do not increment counter
{\small\begin{longtable}[]{@{}
  >{\raggedright\arraybackslash}p{(\linewidth - 2\tabcolsep) * \real{0.5000}}
  >{\raggedright\arraybackslash}p{(\linewidth - 2\tabcolsep) * \real{0.5000}}@{}}
\toprule\noalign{}
\begin{minipage}[b]{\linewidth}\raggedright
Mondo
\end{minipage} & \begin{minipage}[b]{\linewidth}\raggedright
Specializzazione
\end{minipage} \\
\midrule\noalign{}
\endhead
\bottomrule\noalign{}
\endlastfoot
M1 & implementazioni in giurisdizioni \(\Phi\)-engine (Brazil,
Indonesia, alcuni UE) \\
M2 & implementazioni in giurisdizioni \emph{low-blast-radius} (Svizzera,
paesi nordici, Singapore) \\
M3 & implementazioni in giurisdizioni \emph{anti-fragili al decoupling}
(UE, Brazil, Israele) \\
\end{longtable}}
}

\Proved{} \textbf{Condizione antifragile.} La ridondanza modulare è il
principio structural di Taleb (\emph{Antifragile} 2012, capp. 16-17):
paga overhead nello scenario base, produce sopravvivenza nel cigno nero.
Operazionalizzata in \(\mathbb{V}^{(4)}\): aumenta
\(A^{\text{organizzativo}}\) a costo di \(\sim 15\)-\(25\%\) di
efficienza nominale.

\subsection{O-S --- Missione Orientata ai Mondi-Target (orizzonte 3-15
anni)}\label{o-s-missione-orientata-ai-mondi-target-orizzonte-3-15-anni}

\Hypoth{} \textbf{Protocollo.} L'organizzazione dichiara esplicitamente
\emph{quali} dei quattro mondi M1-M4 considera target prioritario, e
calibra missione, comunicazione, alleanze coerentemente.

{\def\LTcaptype{none} % do not increment counter
{\small\begin{longtable}[]{@{}
  >{\raggedright\arraybackslash}p{(\linewidth - 2\tabcolsep) * \real{0.5000}}
  >{\raggedright\arraybackslash}p{(\linewidth - 2\tabcolsep) * \real{0.5000}}@{}}
\toprule\noalign{}
\begin{minipage}[b]{\linewidth}\raggedright
Tipo di organizzazione
\end{minipage} & \begin{minipage}[b]{\linewidth}\raggedright
Target tipico
\end{minipage} \\
\midrule\noalign{}
\endhead
\bottomrule\noalign{}
\endlastfoot
Università, think tank, ONG transnazionale & M1 (ottimale per produzione
di \(\Phi\)-engine) \\
Impresa manifatturiera, farmaceutica, energetica & M3 (riassetto
sistemico veloce, opportunità di riposizionamento) \\
Comunità locali, mutue, cooperative & M2 (resilience al collapse del
centro coerente) \\
Banche centrali, fondi sovrani & M4 (preservare ottica di stabilità
plurale) \\
\end{longtable}}
}

\Proved{} \textbf{Condizione antifragile.} La dichiarazione esplicita
riduce \(\varepsilon_{\text{dis}}^{\text{organizzativo}}\) (gli
stakeholder sanno cosa l'organizzazione perfollows) e libera
\emph{opzionalità} --- l'organizzazione si rende capace di crescere nei
mondi compatibili con la sua missione instead di perseguire
massimizzazione miope nel mondo attuale.



\section{Strategie Statali (livello
S)}\label{strategie-statali-livello-s}

\subsection{S-T --- Protocollo di Coerenza in Crisi (orizzonte
mesi)}\label{s-t-protocollo-di-coerenza-in-crisi-orizzonte-mesi}

\Hypoth{} \textbf{Protocollo.} Lo Stato mantiene, first della crisi, un
\emph{protocollo dichiarato pubblicamente} di risposta coordinata
istituzionale-comunicativa-economica per almeno tre classi di shock:
shock economico-finanziario, shock geopolitico, shock
pandemico/climatico. Il protocollo specifica chi decide cosa, in quanto
tempo, con quale comunicazione.

\Proved{} \textbf{Condizione antifragile.} L'esistenza pubblica del
protocollo \emph{aumenta \(s_{iii}\)} (legitimacy) anche in the absence of
crisi (segnale di competenza istituzionale), e \emph{riduce
\(\varepsilon_{\text{dis}}\)} in caso di crisi attivandosi (la
cittadinanza vede una risposta coordinata anziché frammentata).

\subsection{S-O --- Politica Estera Multi-Vettoriale (orizzonte 1-3
anni)}\label{s-o-politica-estera-multi-vettoriale-orizzonte-1-3-anni}

\Hypoth{} \textbf{Protocollo.} Lo Stato mantiene attivamente
\emph{almeno cinque} canali bilaterali significativi e non sovrapposti:
USA, China, UE, Russia, attore regionale autonomo (per il Brazil: India,
Africa, Mercosul; per uno Stato africano: BRICS, UE, USA, China, Lega
Araba; ecc.). L'\emph{Estado Logístico} di Cervo è il modello di
reference (Ch.~9.1).

{\def\LTcaptype{none} % do not increment counter
{\small\begin{longtable}[]{@{}
  >{\raggedright\arraybackslash}p{(\linewidth - 2\tabcolsep) * \real{0.5000}}
  >{\raggedright\arraybackslash}p{(\linewidth - 2\tabcolsep) * \real{0.5000}}@{}}
\toprule\noalign{}
\begin{minipage}[b]{\linewidth}\raggedright
Mondo
\end{minipage} & \begin{minipage}[b]{\linewidth}\raggedright
Specializzazione
\end{minipage} \\
\midrule\noalign{}
\endhead
\bottomrule\noalign{}
\endlastfoot
M1 & priorità a co-fondazione di iniziative \(\Phi\)-engine (TFFF, BRICS
BMG, Group of Friends for Peace estesi) \\
M2 & priorità alla preservazione delle relazioni minime indipendenti dai
blocchi \\
M3 & priorità a riposizionamento veloce nei nuovi assi post-rottura \\
\end{longtable}}
}

\Proved{} \textbf{Condizione antifragile.} Il multi-vettorialismo è
strutturalmente \(A^{\text{stato}} > 0\) with respect to strategia di
alleanza unilaterale: la perdita di un canale non destabilizza il
sistema diplomatico.

\subsection{\texorpdfstring{S-S --- Investimento Strutturale in
\(\Phi\)-engine (orizzonte 3-15
anni)}{S-S --- Investimento Strutturale in \textbackslash Phi-engine (orizzonte 3-15 anni)}}\label{s-s-investimento-structural-in-phi-engine-orizzonte-3-15-anni}

\Hypoth{} \textbf{Protocollo.} Lo Stato investe in modo programmatico e
quantificato in:

\begin{enumerate}
\def\labelenumi{\arabic{enumi}.}
\tightlist
\item
  \textbf{Infrastruttura comunicativa pluralista} (eredità Topografia
  \S{}5.4: aumenta \(K\) verso \(K_c\));
\item
  \textbf{Educazione su tutti e tre gli assi} \((x, y, z)\) della
  Topografia (aumenta \(Q_0\) collettivo);
\item
  \textbf{Indipendenza giudiziaria + neutralità forze armate} (aumenta
  \(R_c\) --- Ch.~9.4);
\item
  \textbf{Diversificazione produttiva su almeno tre vettori
  non-correlati} (aumenta \(A^{\text{nazionale}}\));
\item
  \textbf{Politica fiscale anti-involutiva} (riduce
  \(\Omega^{\text{strutturale}}\));
\item
  \textbf{Trasparenza della comunicazione governativa} (riduce
  \(\varepsilon_{\text{dis}}^{\text{collettivo}}\)).
\end{enumerate}

\Proved{} \textbf{Condizione antifragile.} Sei interventi su tutte e
quattro le variables di \(\mathbb{V}^{(4)}\) + una variable esogena
(\(K\)) + una eredità topografica (\(Q_0\)). La coerenza dell'insieme è
verificationta: tutti i sei interventi rispettano R1-R6 simultaneamente.

\Hypoth{} \textbf{Tempo caratteristico.} Un investimento S-S coerente
sposta \(\boldsymbol{p}^{\text{stato}}(t)\) verso \(\mathcal{R}_+\) in
\(\sim 5\)-\(8\) anni --- \emph{non} in un ciclo elettorale. Da qui la
difficoltà politica structural: il sistema democratico premia
traiettorie a \(< 4\) anni, while l'intervento \(\Phi\)-engine richiede
\(5\)-\(8\). La risoluzione di questa tensione è il problema
fondamentale della politica nel periodo 2026-2031.



\section{Il Caso Particolare
Brasiliano}\label{il-caso-particolare-Brazilian}

\Hypoth{} \textbf{Premessa.} B1 (esito 4 ottobre 2026) è la predizione
più informativa dell'apparato (Ch.~11.3). Le strategie qui distinguono
\emph{pre-elettorale} (aprile-settembre 2026) e \emph{post-elettorale}
(ottobre 2026 in then).

\subsection{\texorpdfstring{Strategie Pre-elettorali (per massimizzare
\(\Pr(B_1 \text{ vera}))\)}{Strategie Pre-elettorali (per massimizzare \textbackslash Pr(B\_1 \textbackslash text\{ vera\}))}}\label{strategie-pre-elettorali-per-massimizzare-prb_1-text-vera}

\Hypoth{} \textbf{Protocollo per la coalizione Lula/sinistra.} Ridurre
il gap approval/disapproval entro settembre 2026 through azioni che
agiscono \emph{simultaneamente} su quattro variables:

{\def\LTcaptype{none} % do not increment counter
{\small\begin{longtable}[]{@{}
  >{\raggedright\arraybackslash}p{(\linewidth - 4\tabcolsep) * \real{0.3333}}
  >{\raggedright\arraybackslash}p{(\linewidth - 4\tabcolsep) * \real{0.3333}}
  >{\raggedright\arraybackslash}p{(\linewidth - 4\tabcolsep) * \real{0.3333}}@{}}
\toprule\noalign{}
\begin{minipage}[b]{\linewidth}\raggedright
Azione
\end{minipage} & \begin{minipage}[b]{\linewidth}\raggedright
Variabile colpita
\end{minipage} & \begin{minipage}[b]{\linewidth}\raggedright
Tempo richiesto
\end{minipage} \\
\midrule\noalign{}
\endhead
\bottomrule\noalign{}
\endlastfoot
Annuncio + esecuzione di un programma fiscale credibile (riforma
frammentaria già in corso) & \(\Omega \downarrow\) & 4-6 mesi \\
Comunicazione pubblica con \(z\) alto (dialectical complexity, non \(z\)
basso egemonico) & \(Q_0 \uparrow\), \(\Phi^{\text{interno}} \uparrow\)
& 3-6 mesi \\
Esempi pubblici di \(R_c\) effettivo (oltre Bolsonaro: nuove condanne,
regulation tech, environmental enforcement) & \(s_{iii} \uparrow\),
\(\varepsilon_{\text{dis}} \downarrow\) & continuo \\
Successione visibile (Haddad o equivalente come VP designato) &
\(A^{\text{politico}} \uparrow\) (opzione di continuità) & aprile-luglio
2026 \\
\end{longtable}}
}

\Conj{} \textbf{Caveat.} Le strategie sopra sono \emph{operative} nel
senso DEQ$^2$, non normative nel senso politico. La DEQ$^2$ non prescrive
\emph{quale} coalizione debba vincere --- descrive le condizioni
strutturali sotto cui \(T_s^{(2),\text{sys},\text{Brasile}} > 1\) è
preservato. Le ricadute politiche dipendono dal sistema valoriale del
lettore.

\subsection{Strategie Post-elettorali --- Ramo
A/B}\label{strategie-post-elettorali-ramo-ab}

\Hypoth{} \textbf{Protocollo.} Consolidamento operational del
\(\Phi\)-engine: attuazione di S-S (\S{}12.4.3) con priorità a
infrastruttura comunicativa pluralista (1), educazione \(z\)-intensiva
(2), e indipendenza giudiziaria già verificationta (3). Espansione del Group
of Friends for Peace + completamento operational TFFF + lancio operational
BRICS BMG (predizione B2). Articolazione esplicita della \emph{posizione
\(\Phi\)-engine} nei consessi multilaterali --- non lasciarla implicita.

\Proved{} \textbf{Condizione antifragile.} In ramo A/B, M1 si realizza
con \(\Pr \approx 0{,}55\)-\(0{,}70\) --- il protocollo S-S è ottimale
per M1.

\subsection{Strategie Post-elettorali --- Ramo
C}\label{strategie-post-elettorali-ramo-c}

\Hypoth{} \textbf{Protocollo defensivo.} L'esistenza del \emph{$\Phi$-engine}
non dipende solo dall'esecutivo. Le istituzioni \(R_c\) verificationte del
Ch.~9.4 (STF, TSE, civil society, media indipendenti, neutralità forze
armate) devono essere \emph{protette} sopra ogni altra priorità --- sono
la garanzia che lo scenario \emph{worst-C} non si realizzi. Strategie
concrete:

\begin{itemize}
\tightlist
\item
  difesa giuridica preventiva dei magistrati di STF e TSE;
\item
  codificazione delle protezioni per civil servants tecnici (analogo
  Brazilian della \emph{Schedule F} USA che Trump 2.0 ha smantellato
  --- Ch.~7.2);
\item
  diversificazione delle fonti di finanziamento dei media indipendenti;
\item
  mantenimento della neutralità delle forze armate via formazione e
  clausole costituzionali.
\end{itemize}

\Proved{} \textbf{Condizione antifragile.} Anche se M1 fallisce e si
realizza M2, la presenza di \(R_c\) effettivo riduce la profondità della
transizione del Brazil da \(\mathcal{R}_+\) a \(\mathcal{R}_-\). Il
payoff è \emph{positivo} in \emph{tutti} i sotto-scenari del ramo C.



\section{\texorpdfstring{Strategie di Sopravvivenza dentro
\(\mathcal{R}_-\)}{Strategie di Sopravvivenza dentro \textbackslash mathcal\{R\}\_-}}\label{strategie-di-sopravvivenza-dentro-mathcalr_-}

\Hypoth{} \textbf{Premessa onesta.} Cittadini, organizzazioni e attori
sub-statali dentro un sistema in \(\mathcal{R}_-\) (USA è il caso
paradigmatico) non possono attendere il recupero sistemico --- la
diagnosi del Ch.~7 dà \(T_s^{(2),\text{sys},\text{USA}} < 1\) per
l'intero orizzonte 2026-2031 con \(\Pr \approx 0{,}80\). Le loro
strategie devono essere progettate per \emph{prosperare entro la
decoherence}, non per \emph{attenderne la fine}.

\Hypoth{} \textbf{Protocollo S\(_{\mathcal{R}_-}\).} Tre principi
operativi:

\begin{enumerate}
\def\labelenumi{\arabic{enumi}.}
\tightlist
\item
  \textbf{Discoupling parziale dalla traiettoria sistemica}:
  ridurre l'esposizione a indicatori macro (politica fiscale, currency,
  polarizzazione mediatica) through scelte di vita locali, valuta
  diversificata, comunità di affidamento.
\item
  \textbf{Costruzione di sotto-sistemi \(\Phi\)-coerenti}: operare a
  scala where \(\Phi\) può essere ricostruita per scelta --- comunità
  professionali, città medie con governance funzionante, network
  transnazionali specializzati. Sono \emph{isole di coerenza} dentro
  l'oceano decoerente.
\item
  \textbf{Mantenimento dell'opzionalità sovranazionale}: cittadinanza,
  residenza, asset, professione che permettano il riposizionamento entro
  \(5\) anni se la dinamica peggiora oltre lo scenario mediano (verso
  \emph{worst}).
\end{enumerate}

\Proved{} \textbf{Condizione antifragile.} La strategia
\(S_{\mathcal{R}_-}\) è strettamente più robusta di qualsiasi strategia
\emph{all-in} sul recupero USA: paga \(\sim 10\)-\(20\%\) di payoff
atteso nello scenario \emph{best}, ma evita perdite catastrofiche in
\emph{median} e \emph{worst}.



\section{Chapter Summary}\label{sintesi-del-chapter}

\Proved{} \textbf{In una frase:} la DEQ$^2$ non termina nella diagnosi ma
offre ventisette protocolli antifragili (3 livelli $\times$ 3 orizzonti $\times$ 3
scenari significativi), specializzati per ciascuno dei mondi possibili
2031, con condizione antifragile esplicita (payoff positivo in
\(\geq 3\) dei 4 mondi); la strategia Brazilian è scorporata data
l'importanza di B1, e una strategia di sopravvivenza dedicata
\(S_{\mathcal{R}_-}\) è offerta agli attori dentro a sistemi in
phase transition.

{\def\LTcaptype{none} % do not increment counter
{\small\begin{longtable}[]{@{}
  >{\raggedright\arraybackslash}p{(\linewidth - 6\tabcolsep) * \real{0.2500}}
  >{\raggedright\arraybackslash}p{(\linewidth - 6\tabcolsep) * \real{0.2500}}
  >{\raggedright\arraybackslash}p{(\linewidth - 6\tabcolsep) * \real{0.2500}}
  >{\raggedright\arraybackslash}p{(\linewidth - 6\tabcolsep) * \real{0.2500}}@{}}
\toprule\noalign{}
\begin{minipage}[b]{\linewidth}\raggedright
Livello
\end{minipage} & \begin{minipage}[b]{\linewidth}\raggedright
Tattico
\end{minipage} & \begin{minipage}[b]{\linewidth}\raggedright
Operativo
\end{minipage} & \begin{minipage}[b]{\linewidth}\raggedright
Strategico
\end{minipage} \\
\midrule\noalign{}
\endhead
\bottomrule\noalign{}
\endlastfoot
Individuale & I-T opzionalità minima & I-O portafoglio multi-vettoriale
& I-S traiettoria anti-eteronomica \\
Organizzativo & O-T redistribuzione decisionale & O-O ridondanza
modulare & O-S missione orientata ai mondi-target \\
Statale & S-T protocollo di coerenza & S-O politica estera
multi-vettoriale & S-S investimento structural \(\Phi\)-engine \\
\end{longtable}}
}

\Hypoth{} \textbf{Tesi politica conclusiva.} L'opposizione tradizionale
fra \emph{politica realista} (massimizzare in scenario fissato) e
\emph{politica utopica} (ignorare scenario fissato) viene superata dalla
DEQ$^2$: la \emph{politica antifragile} prende sul serio lo spazio dei
mondi possibili e progetta per sopravvivere e prosperare through
essi. La differenza è structural, non retorica.



\section{The Evolutionary Foundation of Strategies: Ostrom and the Commons of \(\Phi\)}\label{il-fondamento-evolutivo-delle-strategie}

\Proved{} \textbf{The empirical (not normative) basis of cooperation.}
The strategies in Ch.~12 prescribe actions that increase $\Phi$ and
reduce $\Omega$. But why are these actions \emph{possible} --- and not
merely desirable? The answer is provided by Elinor Ostrom (1990): the
\emph{commons} --- collective resources --- \emph{can} be managed
cooperatively over the long run when precise institutional conditions are
satisfied. Cooperation is neither inevitable nor utopian: it is
\emph{conditionally stable} under well-designed institutions.

This is the correct normative foundation: not Kropotkin (``cooperation
is in nature,'' which incurs the naturalistic fallacy) but Ostrom
(``cooperation is empirically documented under these conditions'').

\Hypoth{} \textbf{Ostrom's eight principles as conditions for
\(\Phi > 0\).} Ostrom (1990) identifies eight design principles for the
sustainable management of commons. In the DEQ$^2$, these principles
translate directly into conditions for $\Phi > 0$ and $\phi < \phi^*$:

{\def\LTcaptype{none}
{\small\begin{longtable}[]{@{}
  >{\raggedright\arraybackslash}p{(\linewidth - 4\tabcolsep) * \real{0.5000}}
  >{\raggedright\arraybackslash}p{(\linewidth - 4\tabcolsep) * \real{0.5000}}@{}}
\toprule\noalign{}
\begin{minipage}[b]{\linewidth}\raggedright
Ostrom principle (1990)
\end{minipage} & \begin{minipage}[b]{\linewidth}\raggedright
CPGG/DEQ$^2$ correspondence (\Hypoth{})
\end{minipage} \\
\midrule\noalign{}
\endhead
\bottomrule\noalign{}
\endlastfoot
1. Clear boundaries & Reduces $\phi$ by preventing extractive infiltration \\
2. Context-adapted rules & Stabilizes $\sigma_{\text{san}}$ locally \\
3. Participation in rule-making & Increases $\kappa_C$ (cooperator cohesion) \\
4. Monitoring & Reduces $\kappa_E$ (extractor centralization) \\
5. Graduated sanctions & Maintains $\sigma_{\text{san}} > \sigma_{\text{crit}}$ \\
6. Conflict resolution mechanisms & Reduces $\varepsilon_{\text{dis}}$ \\
7. External recognition & Increases $s_{iii}$ (ethical-discursive legitimacy) \\
8. Nested governance & Produces $\Phi$ at multiple levels (scalability) \\
\end{longtable}}
}

\Proved{} \textbf{Interpretation.} The twenty-seven protocols of
Ch.~12.2--12.4 are contextual implementations of Ostrom's eight
principles within the geopolitical reality of 2026--2031. They are not
ideal prescriptions: they are operational protocols for maintaining
$\phi < \phi^*$ while the global system navigates toward the 2028--2031
bifurcation.

\Hypoth{} \textbf{The limit: the scale problem.} Ostrom's commons are
local (lakes, forests, irrigation systems). Extension to geopolitical
scale requires the hypothesis that the design principles scale --- a
hypothesis supported by partial evidence (BRICS+ as an attempt at an
energy/financial commons at regional scale) but not yet demonstrated at
global scale. This is the true stake of the 2026--2031 period: the first
empirical test of the scalability of Ostrom's commons to hegemonic scale.



\section{Closing Epistemic Notes
Chapter}\label{note-epistemiche-di-chiusura-chapter}

{\def\LTcaptype{none} % do not increment counter
{\small\begin{longtable}[]{@{}
  >{\raggedright\arraybackslash}p{(\linewidth - 4\tabcolsep) * \real{0.3333}}
  >{\raggedright\arraybackslash}p{(\linewidth - 4\tabcolsep) * \real{0.3333}}
  >{\raggedright\arraybackslash}p{(\linewidth - 4\tabcolsep) * \real{0.3333}}@{}}
\toprule\noalign{}
\begin{minipage}[b]{\linewidth}\raggedright
\#
\end{minipage} & \begin{minipage}[b]{\linewidth}\raggedright
Claim
\end{minipage} & \begin{minipage}[b]{\linewidth}\raggedright
Status
\end{minipage} \\
\midrule\noalign{}
\endhead
\bottomrule\noalign{}
\endlastfoot
1 & Distinzione strategie passive/antifragili in linguaggio DEQ$^2$ &
\Hypoth{} formalizzazione di Taleb 2012 \\
2 & Matrice 3$\times$3 livelli-orizzonti & \Hypoth{} schema operational \\
3 & Tre protocolli individuali (I-T, I-O, I-S) & \Hypoth{} protocolli
operativi, non normativi \\
4 & Tre protocolli organizzativi (O-T, O-O, O-S) & \Hypoth{} protocolli
operativi \\
5 & Tre protocolli statali (S-T, S-O, S-S) & \Hypoth{} protocolli
operativi \\
6 & Tempo caratteristico S-S = 5-8 anni vs ciclo elettorale 4 anni &
\Hypoth{} problema structural democratico identificato \\
7 & Strategie pre/post-elettorali brasiliane (operative, non normative)
& \Hypoth{} esplicito caveat metodologico \\
8 & Protocollo \(S_{\mathcal{R}_-}\) per cittadini USA-like & \Hypoth{}
onestà sul fatto che la teoria descrive realtà sgradevoli \\
9 & Condizione antifragile esplicita per ogni protocollo & \Hypoth{}
verification formal \\
10 & Tesi conclusiva: politica antifragile vs realista/utopica &
\Hypoth{} posizionamento concettuale \\
\end{longtable}}
}


